% ANCHOR: after_label def:c20-setting
\begin{remark}[Physics notation: the regulated propagator and the finite counterterm]
Every propagator carries the factor $e^{-\varepsilon|k|^2}/|k|^2$: the proper-time cutoff $t\ge\varepsilon$ on the Schwinger representation $1/k^2=\int_0^\infty e^{-tk^2}dt$, a momentum cutoff at $|k|\sim\varepsilon^{-1/2}$ that is smooth and gauge covariant. In position space $D_\varepsilon(r)$ is the massless propagator $1/(4\pi^2r^2)$ smoothed at $r\lesssim\sqrt\varepsilon$. The finite counterterm $\beta_2\int\mathrm{dvol}\,(dA|_{W_{\mathbb R}^\perp},dA|_{W_{\mathbb R}^\perp})=\tfrac{\beta_2}4\int d^4x\,(\partial_mA^{ij},\partial_mA_{ij})$ is a finite scalar wave-function renormalisation $Z_\phi=1+2\beta_2\cdot\frac{g^2N}{32\pi^2}$ in the normalisation $A_I=-igA^a_IT^a$; its value is not fixed by the scalar two-point functions alone and is the scheme parameter carried through this chapter. The operators $O_{ab}$ are the gauge-invariant scalar bilinears $a_Ib_J\operatorname{tr}(\phi_I\phi_J)$ of $\mathcal N=4$ SYM; $K$ is the Konishi operator $\operatorname{tr}(\phi_I\phi_I)$ and the traceless ones form the $\mathbf{20'}$ of $SU(4)_R$, the chiral primaries of dimension two.
\end{remark}

% ANCHOR: after_label prop:c20-selfenergy
\begin{remark}[Physics notation: the scalar self-energy in Feynman gauge]
In the canonical normalisation of the key, the one-loop scalar self-energy $\Sigma(p)$ of $\mathcal N=4$ SYM in Feynman gauge is the sum of the gluon--scalar bubble, the two seagulls (gluon loop and scalar loop) and the gaugino loop; the ghosts do not couple to the scalars. With the proper-time cutoff the quadratically divergent parts $\propto\varepsilon^{-1}$ are the scalar mass terms that the counterterm removes, and the logarithmic part $2\frac{g^2N}{16\pi^2}p^2\log\varepsilon$ is the scalar wave-function renormalisation, $\gamma_\phi=\frac{g^2N}{8\pi^2}$ in Feynman gauge (the block after Section 6.4). The finite part $\frac{g^2N}{16\pi^2}p^2(2\log p^2+2\gamma_E+2\log2-\tfrac{43}8)$ is scheme dependent; its $\log p^2$ carries the anomalous dimension, and the constant is specific to the heat-kernel cutoff. The gaugino loop is the loop of four Weyl fermions $\lambda^i_\alpha$ with the Yukawa couplings $g\,\bar\lambda_{\dot\alpha i}\bar\Sigma^{Iij}[\phi_I,\bar\lambda^{\dot\alpha}{}_j]$-type fixed by supersymmetry, equivalently the ten-dimensional Majorana--Weyl loop with external polarisations in the internal directions.
\end{remark}

% ANCHOR: after_label thm:c20-result
\begin{remark}[Physics notation: the two-point functions and the Konishi anomalous dimension]
Equation \eqref{eq:c20-result} is the statement that the two-point function of the $\mathbf{20'}$ chiral primaries is $\frac{C}{|x|^4}\big(1-\frac{g^2N}{32\pi^2}(\tfrac{27}2+4\log\tfrac32-8\beta_2)\big)$ with no logarithm (protected dimension $2$), and that the Konishi two-point function is $\frac{C_K}{|x|^{4+2\gamma_K}}$ with $\gamma_K=\frac{3g^2N}{4\pi^2}$, the coefficient $24$ of $\log|x|^2$ in units of $-\frac{g^2N}{32\pi^2}$ being $-\gamma_K$. The exchange graph (EX) is the gluon exchange between the two scalar lines, the quartic graph (QU) the scalar contact interaction, (SE) the self-energy insertions, (OC) the counterterms of the composite operators (the renormalisation constant $Z_O$), and the identity $(k+k')\cdot(2q-k-k')=2q^2-\dots$ is the standard trick that turns the gluon exchange between scalar lines into bubbles by cancelling propagators against numerators; with the heat-kernel cutoff each cancelled propagator leaves a Gaussian contact of range $\sqrt\varepsilon$, whose localisation produces the $\log(2/3)$ terms. The one-loop constants of the $\mathbf{20'}$ and of Konishi are scheme data: they change under finite wave-function renormalisation ($\beta_2$), finite operator renormalisation and the choice of the unit of length.
\end{remark}

% ANCHOR: after_label prop:c20-gram
\begin{remark}[Physics notation: radial quantisation and the Zamolodchikov metric]
The radial inner product of Definition \ref{def:c8-inner} is the inner product of radial quantisation: the state $|O\rangle=O(0)|0\rangle$ and its conjugate $\langle O|=\lim_{x\to\infty}|x|^{2\Delta}\langle0|O^\dagger(x)$, so that $\langle O_1|O_2\rangle$ is the two-point coefficient at unit separation, the Zamolodchikov metric on the space of operators. The inversion $x\mapsto x/x^2$ composed with complex conjugation of the fields realises the conjugation $O\mapsto O^\dagger$ of the CFT; for the bilinears it conjugates the $SO(6)$ polarisations $a,b$. Equation \eqref{eq:c20-gram} is this metric at one loop on the twenty-one scalar bilinears: the $\mathbf{20'}$ block is proportional to the tree metric (the protected normalisation up to the scheme constant), the Konishi norm carries in addition $24(\gamma_E-3\log2)$ units, the finite part that accompanies the anomalous dimension in this scheme. Hermiticity and positivity of the metric are the unitarity of the one-loop theory on this sector.
\end{remark}

% ANCHOR: after_label thm:c20-hermrad
\begin{remark}[Physics notation: reflection positivity and Hermiticity]
The reflection $\Theta_n$ is the Osterwalder--Schrader reflection in the hyperplane orthogonal to $n$; the reflection form $\langle a^{\Theta_n},b\rangle$ is the Euclidean realisation of the Hilbert-space inner product, and its Hermiticity is the statement that the Euclidean path integral with a real action and a reflection-symmetric regulator defines a Hermitian form. The proposition shows that this survives the heat-kernel cutoff, the counterterms and the finite counterterms whenever they are reflection symmetric, to every order; a non-real finite counterterm (a complex $\beta_2$, a complex wave-function renormalisation) is exactly what would break it. The radial form is the reflection form transported by the conformal map from the half space to the ball, generated by $P_n-K_n$; at one loop the transport uses the one-loop translation and special conformal generators, and the Ward identity along $P_n-K_n$ is what makes the transport unitary.
\end{remark}
